\documentclass[11pt,a4paper]{article}
\usepackage{fontspec}
\usepackage{polyglossia}
\setdefaultlanguage{russian}
\setotherlanguage{english}
\setmainfont{Times New Roman}
\setsansfont{Arial}
\setmonofont{Consolas}
\usepackage[margin=2.2cm]{geometry}
\usepackage{amsmath,amssymb,amsthm,mathtools}
\usepackage{booktabs}
\usepackage{enumitem}
\usepackage[colorlinks=true,linkcolor=blue!60!black,citecolor=blue!60!black,urlcolor=blue!60!black]{hyperref}

\newtheorem{theorem}{Теорема}[section]
\newtheorem{corollary}[theorem]{Следствие}
\newtheorem{proposition}[theorem]{Предложение}
\theoremstyle{definition}
\newtheorem{definition}[theorem]{Определение}
\theoremstyle{remark}
\newtheorem{remark}[theorem]{Замечание}

\title{\textbf{Динамика траекторий Коллатца на естественных модулярных кольцах $\mathbb{Z}_{2^a}$ и $\mathbb{Z}_{2^a - 3^b}$:\\ разложение показателей в базисе $\mathrm{SL}_2(\mathbb{Z})$ Каталана, двусторонние границы дрейфа и знакоположительные суммы Вронскиана в Lean~4}}
\author{Дмитрий Каратаев \and Луч (AI-Guiders)}
\date{2026}

\begin{document}
\maketitle

\begin{abstract}
Задача Коллатца $3n+1$ переформулирована как задача о дискретном движении на двусортной алгебраической системе носителей $\langle\mathcal{M},\mathcal{T}\rangle$ (в смысле А.\,И.~Мальцева) над естественными модулярными кольцами $\mathbb{Z}_{2^a} = \mathbb{Z}_{4Q}$ ($Q = 2^{a-2}$) и кольцом зазора Пиллаи $\mathbb{Z}_{C_{a,b}}$ ($C_{a,b} = 2^a - 3^b$). Во-первых, обобщая кусочно-линейный ротатор и стереографическое исчисление Вейерштрасса на произвольное четырёхквадрантное кольцо $\mathbb{Z}_{4Q}$ (\texttt{thm18} в \texttt{RotatorCalculusPaper.lean}), мы доказываем, что каждая $a$-шаговая бинарно-троичная траектория $2^a n_{\mathrm{final}} - 3^b n_{\mathrm{start}} = D_{a,b}$ раскладывается в точный баланс частного и остатка $2^a(n_{\mathrm{final}} - n_{\mathrm{start}}) + (2^a - 3^b)n_{\mathrm{start}} \equiv D_{a,b}$, где полином дрейфа $D_{a,b} = \sum_{j=0}^{b-1} 3^{b-1-j} 2^{s_j}$ безусловно зажат в геометрической вилке Пиллаи $3^b - 2^b \le D_{a,b} \le 2^{a-b}(3^b - 2^b)$, а $a$-битный вычет $r = n_{\mathrm{start}}\bmod 2^a\in\mathbb{Z}_{2^a}$ биективно определяет вектор чётностей (сжимая траекторию на $382/512 = 74{,}61\%$ вычетов при $a=9$). Во-вторых, для любых двух последовательных дуг траектории $P_1 n_{\mathrm{mid}} - Q_1 n_{\mathrm{start}} = D_1$ и $P_2 n_{\mathrm{final}} - Q_2 n_{\mathrm{mid}} = D_2$ и суммарный дрейф $D_{\mathrm{tot}} = Q_2 D_1 + P_1 D_2$, и суммарный зазор Пиллаи $C_{\mathrm{tot}} = Q_2 C_1 + P_1 C_2$ порождаются одним и тем же линейным оператором $(Q_2, P_1)$, связывая хордовое смещение на любом замкнутом цикле ($n_{\mathrm{final}} = n_{\mathrm{start}}$) с определителем Вронскиана $2\times 2$: $C_{\mathrm{tot}}(n_{\mathrm{mid}} - n_{\mathrm{start}}) \equiv C_2 D_1 - C_1 D_2$. В-третьих, доказано, что унимодулярная матрица $\mathcal{B}_{2,3} = \begin{pmatrix} 2 & 3 \\ 1 & 2 \end{pmatrix} \in \mathrm{SL}_2(\mathbb{Z})$ единственным образом раскладывает любую пару показателей $(a,b)\in\mathbb{Z}^2$ по базису Каталана $u = 2a - 3b$ (дуга $(2,1)$ с зазором $2^2 - 3^1 = +1$) и $v = 2b - a$ (дуга Каталана $(3,2)$ с зазором $3^2 - 2^3 = +1$), так что $2^a = 4^u 8^v$ и $3^b = 3^u 9^v$ при $u \ge 1, v \ge 1$ на каждом сжимающем цикле ($\frac{3}{2} < \frac{a}{b} < 2$). Разрезание любого замкнутого цикла на границе между блоком $(4^u,3^u)$ ($G_1 = 4^u - 3^u > 0$) и блоком $(8^v,9^v)$ ($G_2 = 9^v - 8^v > 0$) превращает хордовое уравнение Вронскиана в строго положительную сумму без вычитания: $(2^a - 3^b)(n_{\mathrm{start}} - n_{\mathrm{mid}}) \equiv (9^v - 8^v)D_1 + (4^u - 3^u)D_2 > 0$ (\texttt{thm19}, \texttt{0 sorryAx}, \texttt{0 Classical.choice}).
\end{abstract}

\section{Введение и постановка задачи на естественных кольцах $\mathbb{Z}_{2^a}$ и $\mathbb{Z}_{2^a - 3^b}$}
\label{sec:intro}

Всюду в статье термин \emph{носитель} (\emph{carrier}) употребляется строго в смысле классической теории моделей А.\,И.~Мальцева~\cite{Malcev1970} как основное множество сорта в двусортной алгебраической системе $\langle\mathcal{M},\mathcal{T}\rangle$. В базовой работе по ротаторному исчислению~\cite{KarataevRotator1} дискретный сорт $\mathcal{M}$ был задан на самозамкнутом регистре $\mathbb{Z}_{512}$ ($2^9 = 512$). Фундаментальное свойство построенного аппарата состоит в том, что кусочно-линейный ротатор $(\mathcal{C}_{4Q},\mathcal{S}_{4Q})$, его постоянный разностный вронскиан $\mathcal{W}_{4Q} \equiv Q$ и стереографическое погружение Вейерштрасса переносятся без изменений на \textbf{произвольное} четырёхквадрантное кольцо $\mathbb{Z}_{4Q}$ ($Q\in\mathbb{Z}$, Теорема~\ref{thm:arbitrary-modulus-wronskian}).

Рассмотрим сжатый шаг оператора Коллатца на натуральных числах:
\begin{equation}
  \label{eq:collatz-step}
  T(x) \;=\;
  \begin{cases}
    x / 2, & x \equiv 0 \pmod 2, \\
    (3x + 1) / 2, & x \equiv 1 \pmod 2.
  \end{cases}
\end{equation}
За $a \ge 1$ двоичных делений, содержащих $b \in \{0,\dots,a\}$ нечётных шагов на позициях $0 \le s_0 < s_1 < \dots < s_{b-1} \le a - 1$, траектория из точки $n_{\mathrm{start}}$ в точку $n_{\mathrm{final}} = T^{(a)}(n_{\mathrm{start}})$ удовлетворяет точному линейному диофантову уравнению движения:
\begin{equation}
  \label{eq:collatz-equation-of-motion}
  2^a \cdot n_{\mathrm{final}} \;-\; 3^b \cdot n_{\mathrm{start}} \;=\; D_{a,b}, \qquad
  D_{a,b} \;:=\; \sum_{j=0}^{b-1} 3^{b-1-j}\,2^{s_j}.
\end{equation}
Представляя $n_{\mathrm{start}} = 2^a q + r$ через частное $q = \lfloor n_{\mathrm{start}}/2^a\rfloor \ge 0$ и фазовый остаток $r = n_{\mathrm{start}}\bmod 2^a \in \mathbb{Z}_{2^a}$, замечаем, что вектор чётностей $(e_0,\dots,e_{a-1})\in\{0,1\}^a$ биективно определяется остатком $r\in\mathbb{Z}_{2^a}$ и не зависит от $q$. Вследствие этого число нечётных шагов $b = b(r)$ подчиняется на $\mathbb{Z}_{2^a}$ биномиальному распределению Паскаля $\binom{a}{b}$ (при $a=9$ на $\mathbb{Z}_{512}$ ровно $382/512 = 74{,}61\%$ вычетов имеют $b(r)\le 5$ и $3^{b(r)}\le 243 < 512$), причём
\begin{equation}
  \label{eq:block-affine-transition}
  T^{(a)}(2^a q + r) \;=\; 3^{b(r)}\,q \;+\; x_a(r), \qquad
  2^a\,x_a(r) \;=\; 3^{b(r)}\,r \;+\; D_{a,b}(r),
\end{equation}
а по всем $\binom{a}{b}$ наборам позиций $0 \le s_0 < \dots < s_{b-1} \le a-1$ полином дрейфа $D_{a,b}$ зажат в геометрической вилке Пиллаи:
\begin{equation}
  \label{eq:drift-envelope}
  3^b - 2^b \;=\; \sum_{j=0}^{b-1} 3^{b-1-j}\,2^j \;\le\; D_{a,b} \;\le\; 2^{a-b}\sum_{j=0}^{b-1} 3^{b-1-j}\,2^j \;=\; 2^{a-b}\,\bigl(3^b - 2^b\bigr).
\end{equation}

\section{Ротаторное исчисление на произвольном модуле $\mathbb{Z}_{4Q}$ и закон Вронскиана двух дуг}
\label{sec:two-arc-wronskian}

\begin{theorem}[Погружение ротатора на произвольном модуле $\mathbb{Z}_{4Q}$ и закон Вронскиана двухдуговых траекторий, \texttt{thm18} в \texttt{RotatorCalculusPaper.lean}]
\label{thm:arbitrary-modulus-wronskian}
Для любых целых чисел $Q, A, B, C, u, v, w, \varphi_0, P_1, Q_1, P_2, Q_2, n_{\mathrm{start}}, n_{\mathrm{mid}}, n_{\mathrm{final}}, D_1, D_2 \in \mathbb{Z}$:
\begin{enumerate}[label=\textbf{(\roman*)},leftmargin=*]
\item \textbf{Ротатор и стереографический инвариант на произвольном кольце $\mathbb{Z}_{4Q}$:}
  для $\mathcal{C}_{4Q}(\varphi_0) = Q - \varphi_0$ и $\mathcal{S}_{4Q}(\varphi_0) = \varphi_0$ выполнены $L^1$-норма $\mathcal{C}_{4Q}(\varphi_0) + \mathcal{S}_{4Q}(\varphi_0) \equiv Q$ и постоянство разностного вронскиана $\mathcal{W}_{4Q}(\varphi_0) = \mathcal{C}_{4Q}(\varphi_0)\mathcal{S}_{4Q}(\varphi_0+1) - \mathcal{S}_{4Q}(\varphi_0)\mathcal{C}_{4Q}(\varphi_0+1) \equiv Q$. Для любой тройки $A + B = C$ с парой Безу $A v - B u = 1$ ($u + v = w$) при $I_{\mathrm{rot},Q} = B\mathcal{C}_{4Q}(\varphi_0) - A\mathcal{S}_{4Q}(\varphi_0)$ и $\Delta_{\mathrm{rot},Q} = \mathcal{C}_{4Q}(\varphi_0)v - \mathcal{S}_{4Q}(\varphi_0)u$ тождественно выполняются
  \begin{equation}
    \label{eq:arbitrary-q-invariant}
    Q\,A \;=\; C\,\mathcal{C}_{4Q}(\varphi_0) - I_{\mathrm{rot},Q}, \qquad
    Q\,B \;=\; C\,\mathcal{S}_{4Q}(\varphi_0) + I_{\mathrm{rot},Q}, \qquad
    C\cdot\Delta_{\mathrm{rot},Q} \;-\; w\cdot I_{\mathrm{rot},Q} \;\equiv\; Q;
  \end{equation}
\item \textbf{Композиция двух дуг и разложение смещения в кольце $\mathbb{Z}_{C_{\mathrm{tot}}}$:}
  для двух последовательных дуг траектории $P_1 n_{\mathrm{mid}} - Q_1 n_{\mathrm{start}} = D_1$ и $P_2 n_{\mathrm{final}} - Q_2 n_{\mathrm{mid}} = D_2$ ($P_i = 2^{a_i},\ Q_i = 3^{b_i},\ C_i = P_i - Q_i,\ C_{\mathrm{tot}} = P_1 P_2 - Q_1 Q_2,\ D_{\mathrm{tot}} = Q_2 D_1 + P_1 D_2$) суммарная траектория удовлетворяет
  \begin{equation}
    \label{eq:composite-displacement}
    (P_1 P_2)\,n_{\mathrm{final}} \;-\; (Q_1 Q_2)\,n_{\mathrm{start}} \;=\; D_{\mathrm{tot}}, \qquad
    (P_1 P_2)\,\bigl(n_{\mathrm{final}} - n_{\mathrm{start}}\bigr) \;+\; C_{\mathrm{tot}}\,n_{\mathrm{start}} \;\equiv\; D_{\mathrm{tot}};
  \end{equation}
\item \textbf{Замок Крамера на цикле и хордовое тождество Вронскиана $2\times 2$:}
  на любом замкнутом цикле ($n_{\mathrm{final}} = n_{\mathrm{start}}$) обе вершины $n_{\mathrm{start}}$ и $n_{\mathrm{mid}}$ выражаются через знаменатель Пиллаи $C_{\mathrm{tot}}$, а их хордовое смещение равно определителю Вронскиана $2\times 2$ между зазорами Пиллаи и дрейфами дуг:
  \begin{equation}
    \label{eq:cycle-wronskian-chord}
    C_{\mathrm{tot}}\,n_{\mathrm{start}} = Q_2 D_1 + P_1 D_2, \qquad
    C_{\mathrm{tot}}\,n_{\mathrm{mid}} = P_2 D_1 + Q_1 D_2, \qquad
    C_{\mathrm{tot}}\,\bigl(n_{\mathrm{mid}} - n_{\mathrm{start}}\bigr) \;\equiv\; C_2 D_1 - C_1 D_2.
  \end{equation}
\end{enumerate}
\end{theorem}

\section{Разложение показателей в базисе $\mathrm{SL}_2(\mathbb{Z})$ Каталана и знакоположительная сумма Вронскиана}
\label{sec:catalan-basis}

При произвольном разбиении цикла на две дуги в~\eqref{eq:cycle-wronskian-chord}, если $C_1$ и $C_2$ имеют одинаковый знак, правая часть $C_2 D_1 - C_1 D_2$ содержит вычитание. Покажем, что разложение вектора показателей $(a,b)$ в каноническом базисе $\mathrm{SL}_2(\mathbb{Z})$, порождённом двумя единичными замками Каталана $E_0 = (2,1)$ ($2^2 - 3^1 = +1$) и $E_1 = (3,2)$ ($3^2 - 2^3 = +1$), полностью устраняет вычитание.

\begin{theorem}[Разложение показателей в базисе $\mathrm{SL}_2(\mathbb{Z})$ Каталана и знакоположительная сумма Вронскиана на замкнутых циклах, \texttt{thm19} в \texttt{RotatorCalculusPaper.lean}]
\label{thm:catalan-basis-positive-wronskian}
Для любой пары показателей $(a,b)\in\mathbb{Z}^2$ и любых двух последовательных дуг траектории $P_1 n_{\mathrm{mid}} - Q_1 n_{\mathrm{start}} = D_1$ и $P_2 n_{\mathrm{final}} - Q_2 n_{\mathrm{mid}} = D_2$:
\begin{enumerate}[label=\textbf{(\roman*)},leftmargin=*]
\item \textbf{Унимодулярный базис Каталана $(2,1)$ и $(3,2)$ в $\mathrm{SL}_2(\mathbb{Z})$:}
  матрица базисных показателей $\mathcal{B}_{2,3} = \begin{pmatrix} 2 & 3 \\ 1 & 2 \end{pmatrix}$ имеет определитель $\det\mathcal{B}_{2,3} = 2\times 2 - 3\times 1 = +1$ и единственным образом раскладывает всякую пару $(a,b)\in\mathbb{Z}^2$ по координатам
  \begin{equation}
    \label{eq:catalan-uv-coords}
    u \;:=\; 2a - 3b, \qquad v \;:=\; 2b - a, \qquad
    \begin{pmatrix} a \\ b \end{pmatrix} \;=\; u\begin{pmatrix} 2 \\ 1 \end{pmatrix} + v\begin{pmatrix} 3 \\ 2 \end{pmatrix}, \qquad
    u + v \;=\; a - b.
  \end{equation}
  Для каждого сжимающего цикла Коллатца ($2^a > 3^b \implies \frac{a}{b} > \log_2 3 > \frac{3}{2}$ и $a < 2b \implies \frac{a}{b} < 2$) обе координаты строго положительны ($u \ge 1,\ v \ge 1$), раскладывая полные степени без дробей в
  \begin{equation}
    \label{eq:power-factorization-4-3-8-9}
    2^a \;=\; (2^2)^u\cdot(2^3)^v \;=\; 4^u\cdot 8^v, \qquad
    3^b \;=\; (3^1)^u\cdot(3^2)^v \;=\; 3^u\cdot 9^v;
  \end{equation}
\item \textbf{Разложение зазора Пиллаи $C_{\mathrm{tot}} = 2^a - 3^b$ по единичным замкам $G_1 = 4^u - 3^u > 0$ и $G_2 = 9^v - 8^v > 0$:}
  при $(P_1,Q_1) = (4^u,3^u)$ ($G_1 := P_1 - Q_1 = 4^u - 3^u > 0$, кратно $4 - 3 = 1$) и $(P_2,Q_2) = (8^v,9^v)$ ($G_2 := Q_2 - P_2 = 9^v - 8^v > 0$, кратно $9 - 8 = 1$) полный зазор Пиллаи $C_{\mathrm{tot}} = P_1 P_2 - Q_1 Q_2$ удовлетворяет тождествам
  \begin{equation}
    \label{eq:gap-split-g1-g2}
    C_{\mathrm{tot}} \;\equiv\; 8^v\,\bigl(4^u - 3^u\bigr) \;-\; 3^u\,\bigl(9^v - 8^v\bigr) \;\equiv\; 9^v\,\bigl(4^u - 3^u\bigr) \;-\; 4^u\,\bigl(9^v - 8^v\bigr);
  \end{equation}
\item \textbf{Строго положительная сумма Вронскиана на любом замкнутом цикле ($n_{\mathrm{final}} = n_{\mathrm{start}}$):}
  разрезание любого замкнутого цикла в промежуточной вершине $n_{\mathrm{mid}}$ после первых $2u$ двоичных шагов ($b_1 = u$ нечётных шагов, $P_1 = 4^u,\ Q_1 = 3^u$, дрейф $D_1 \ge 1$; вторая дуга $a_2 = 3v,\ b_2 = 2v$, $P_2 = 8^v,\ Q_2 = 9^v$, дрейф $D_2 \ge 1$) придаёт зазорам противоположные знаки $C_1 = +G_1 > 0$ и $C_2 = -G_2 < 0$, превращая тождество Вронскиана~\eqref{eq:cycle-wronskian-chord} в \textbf{сумму двух строго положительных слагаемых}:
  \begin{equation}
    \label{eq:positive-wronskian-sum}
    \bigl(2^a - 3^b\bigr)\cdot\bigl(n_{\mathrm{start}} - n_{\mathrm{mid}}\bigr) \;\equiv\; \underbrace{\bigl(9^v - 8^v\bigr)\,D_1}_{> 0} \;+\; \underbrace{\bigl(4^u - 3^u\bigr)\,D_2}_{> 0} \;>\; 0,
  \end{equation}
  что безусловно влечёт $n_{\mathrm{start}} > n_{\mathrm{mid}}$ и ограничивает сверху как $(9^v - 8^v)D_1$, так и $(4^u - 3^u)D_2$ величиной $(2^a - 3^b)(n_{\mathrm{start}} - n_{\mathrm{mid}})$.
\end{enumerate}
\end{theorem}

\begin{table}[htbp]
\centering
\caption{Первые шесть резонансных ступеней Штерна--Броко для $\log_2 3$ и их координаты Каталана $(u,v)$.}
\label{tab:collatz-stages}
\small
\begin{tabular}{@{}llllll@{}}
\toprule
\textbf{Ступень} & \textbf{Показатели $(a,b)$} & \textbf{Координаты $(u,v)$} & \textbf{Степени $(2^a, 3^b)$} & \textbf{Зазор $C_{a,b} = 2^a - 3^b$} & \textbf{Линейная комбинация} \\
\midrule
$E_0$ & $(2, 1)$ & $(1, 0)$ & $(4, 3)$ & $+1$ & Тривиальный цикл $n = 1/(4-3) = 1$ \\
$E_1$ & $(3, 2)$ & $(0, 1)$ & $(8, 9)$ & $-1$ & Замок Каталана~I ($3^2 - 2^3 = 1$) \\
$E_2$ & $(5, 3)$ & $(1, 1)$ & $(32, 27)$ & $+5$ & $3^2(+1) + 2^2(-1) = 9 - 4 = +5$ \\
$E_3$ & $(8, 5)$ & $(1, 2)$ & $(256, 243)$ & $+13$ & $3^3(-1) + 2^3(+5) = -27 + 40 = +13$ \\
$E_4$ & $(11, 7)$ & $(1, 3)$ & $(2048, 2187)$ & $-139$ & $3^5(-1) + 2^3(+13) = -243 + 104 = -139$ \\
$E_5$ & $(19, 12)$ & $(2, 5)$ & $(524288, 531441)$ & $-7153$ & $3^7(+13) + 2^8(-139) = -7153$ \\
\bottomrule
\end{tabular}
\end{table}

\section{Формальная верификация в Lean~4}
Теоремы~\ref{thm:arbitrary-modulus-wronskian} и~\ref{thm:catalan-basis-positive-wronskian} (\texttt{thm18} и \texttt{thm19}) формально верифицированы в модуле \texttt{RotatorCalculusPaper.lean} в репозитории \url{https://github.com/AI-Guiders/scientific-papers} (\texttt{0 sorryAx}, \texttt{0 Classical.choice}).

\begin{thebibliography}{9}
\bibitem{Malcev1970}
А.\,И.~Мальцев, \emph{Алгебраические системы}, М.: Наука, 1970 (англ.\ пер.: Springer-Verlag, 1973).
\bibitem{KarataevRotator1}
Д.~Каратаев, Луч, «Ротаторное исчисление на кубооктаэдре: дискретная экспонента Эйлера на $\mathbb{Z}_{512}$, универсальный константный вронскиан и замкнутые квадратуры на кубооктаэдре $O_2(3)$ в Lean~4», \url{https://github.com/AI-Guiders/scientific-papers}, 2026.
\bibitem{Lagarias2010}
J.\,C.~Lagarias (ed.), \emph{The Ultimate Challenge: The $3x+1$ Problem}, American Mathematical Society, Providence, RI, 2010.
\bibitem{Moura2021}
L.~de~Moura and S.~Ullrich, ``The Lean~4 theorem prover and programming language,'' in \emph{CADE~28}, LNCS vol.~12699, Springer, 2021, pp.~625--635.
\end{thebibliography}

\end{document}
